% New finite bound; independent review pending. No master-file edits.
\subsection{An explicit finite bound for the selected later failure}

\begin{result}{C.2.14}{Proved}{A weighted augmentation bound in the five-coordinate group}
\label{lem:ghost-weighted-augmentation}
Let \(p\ge5\) be a prime, let \(a\ge1\), and put \(M=p^a\).
Let \(P_M\) be the free group of nilpotency class at most three
and exponent dividing \(M\) on \(x,y\). It is a finite \(p\)-group
of order \(M^5\). Write
\[
 c=[x,y],\qquad d=[c,x],\qquad e=[c,y],
 \qquad [u,v]=uvu^{-1}v^{-1}.
 \tag{FB.1}
\]
In the group algebra \(\mathbb F_p[P_M]\), its augmentation
ideal \(I\) satisfies
\[
 I^{\,1+10(M-1)}=0.
 \tag{FB.2}
\]
\end{result}
\begin{proof}
The integral Hall model of Lemma~\ref{lem:ghost-hall-model},
reduced modulo \(M\), gives the unique expressions
\(x^i y^j c^h d^u e^v\), with each exponent in
\(\mathbb Z/M\). Its multiplication polynomials have
denominators dividing products of two and three.
These denominators are units modulo \(M\).
Equivalently one can use the truncated Baker--Campbell--Hausdorff
law on the free class-three Lie algebra over \(\mathbb Z/M\);
the two weight-one, one weight-two, and two weight-three
coordinates give \(M^5\) elements.
The powers of any element are the integer multiples of its
logarithm, so the exponent divides \(M\).
The Hall generators and their commutators generate all five
coordinates, which also proves the stated universal property.

Put \(X=x-1,\ Y=y-1,\ C=c-1,\ D=d-1,\ E=e-1\).
The ordered augmentation monomials
\[
 X^iY^jC^hD^uE^v,\qquad 0\le i,j,h,u,v<M,
 \tag{FB.3}
\]
are a basis: expand the powers by the binomial formula in
each variable, in this fixed order. The change from (FB.3)
to the ordered group basis is triangular with diagonal one.
Also every \(M\)-th power of one of these five variables
is zero, since in characteristic \(p\) it is \(g^M-1\).

Assign weights \(1,1,2,3,3\) to \(X,Y,C,D,E\).
The exact group identities following from (FB.1) give
\[
\begin{split}
 YX&=XY+(1+X)(1+Y)
       \bigl((1+C)^{-1}(1+D)^{-1}(1+E)^{-1}-1\bigr),\\
 CX&=XC+(1+X)(1+C)D,\\
 CY&=YC+(1+Y)(1+C)E.
\end{split}
\tag{FB.4}
\]
The variables \(D,E\) commute with all five variables.
All inverse factors in (FB.4) are finite geometric
series, because the corresponding variable has \(M\)-th
power zero. Every summand on the right has weight at
least the weight of the left: respectively two, three,
and three.

Here is a terminating collection procedure, so this
weight assertion applies to products of arbitrary basis
monomials. First count inversions between the \(X\)'s
and \(Y\)'s, ignoring \(C,D,E\); next count occurrences
of a \(C\) preceding an \(X\) or a \(Y\). Move \(D,E\)
to the right. Apply the first rule of (FB.4) to a
\(YX\) inversion, moving intervening \(C\)'s right by
the other two rules when necessary. In every summand
of the first rule the number of \(X,Y\) inversions
decreases; no new \(X\) or \(Y\) is added. In every
summand of the other two rules the first count does
not increase, and the second decreases. A first-rule
step can introduce only finitely many \(C\)'s.
Thus lexicographic induction on these two nonnegative
counts terminates in ordered monomials. Terms with
an exponent at least \(M\) vanish. At each step weight
does not decrease.

Consequently, if \(J_n\) is the span of the basis
monomials of weight at least \(n\), then
\(J_rJ_s\subseteq J_{r+s}\). The augmentation ideal
is \(J_1\), because its quotient by these nonconstant
basis monomials is the one-dimensional augmentation
quotient. The largest weight in (FB.3) is
\((1+1+2+3+3)(M-1)=10(M-1)\).
This proves (FB.2).
\end{proof}

\begin{result}{C.2.15}{Proved}{An empty local fibre at an explicit integral class}
\label{prop:ghost-finite-local-bound}
Let \(p\ge5\), let \(q\ne p\) be a rational prime with
\(q\equiv1\pmod p\), and let \(m\ne0\) be an integer.
Put \(k=v_p(m)\) and
\[
 B(p,m)=10\bigl(p^{k+1}-1\bigr).
 \tag{FB.5}
\]
Use the tangential splitting at zero for
\(U=\mathbb P^1_{\mathbb Q}\setminus\{0,1,\infty\}\).
In the local geometrically pro-\(p\) arithmetic quotient
of integral nilpotency class \(B(p,m)\), there is no
continuous section with both of the following data:
its tame inertia has abelian coordinates \((m,0)\);
and its geometric Frobenius coordinate in the
class-two exponent-\(p\) quotient \(H_p\) is a
nonzero central element.
\end{result}
\begin{proof}
Write \(s=B(p,m)\), \(M=p^{k+1}\), and \(P=P_M\).
The quotient from the free pro-\(p\) geometric group
to \(P\) is characteristic: it is its universal
class-three exponent-\(M\) quotient.
In particular the actual arithmetic action descends
to \(P\). There is also the actual quotient
\(P\to H_p\), obtained by reducing its first three
Hall coordinates modulo \(p\) and killing weight
three.

Suppose that a section with the specified data
exists. The tangential action at \(q\) is unramified
in every integral nilpotency class, as proved in
Theorem~\ref{thm:ghost-fixed-depth-survival}.
Wild inertia and the prime-to-\(p\) part of tame
inertia therefore have zero geometric image.
Let \(T,F\) be tame \(p\)-inertia and Frobenius, with
\(FTF^{-1}=T^q\), and write the geometric coordinates
of the section as \(V,U\), respectively.
Let \(\varphi_q\) be the actual tangential Frobenius
automorphism. Then
\[
 U\varphi_q(V)U^{-1}=V^q,\qquad
 \varphi_q(x)=x^q,\qquad
 \operatorname{gr}_i(\varphi_q)=q^i .
 \tag{FB.6}
\]

We first work only in the rational class-\(s\)
Malcev group. Successively through its central
layers, the nonzero scalars \(q^i-1\) solve
\[
 W^{-1}U\varphi_q(W)=1.
 \tag{FB.7}
\]
This constructs \(W\) in that rational group.
The element \(W^{-1}VW\) then satisfies
\(\varphi_q(W^{-1}VW)=(W^{-1}VW)^q\).
Its abelian image is \((m,0)\).
Compare it with \(x^m\) in successive central
layers: the difference in weight \(i\ge2\) is
killed by \(q^i-q\), a nonzero rational scalar.
Thus
\[
 W^{-1}VW=x^m
 \quad\hbox{in the rational class-\(s\) group}.
 \tag{FB.8}
\]

We use (FB.8) through the integral cyclic-word
quotient, without evaluating \(W\) in a finite
group algebra. The truncated Magnus expansions
of \(V\) and \(x^m\) have coefficients in
\(\mathbb Z_p\). Conjugation in the rational
truncated algebra implies that the sum of their
coefficient differences on each cyclic word is
zero through degree \(s\).
For each degree, the quotient of the free
\(\mathbb Z_p\)-module on words by cyclic rotations
is a free module, with the cyclic orbits as basis.
Therefore these integral coefficient differences
are integral sums of commutators through that
degree.

Now evaluate only these integral expressions at
\(X=x-1,\ Y=y-1\) in \(\mathbb F_p[P]\).
Lemma~\ref{lem:ghost-weighted-augmentation}
gives \(I^{s+1}=0\), so terms of degree greater
than \(s\) vanish. Hence the images of \(V\)
and \(x^m\) have the same image in
\(\mathbb F_p[P]/[\mathbb F_p[P],\mathbb F_p[P]]\).
This quotient has the conjugacy classes of \(P\)
as basis: the relations \(ab-ba\) identify
conjugate group elements, and conversely every
such relation is a difference within one
conjugacy class. Equality of two group basis
elements therefore means they are conjugate.
Choose \(h\in P\) with \(h^{-1}Vh=x^m\).

Conjugating the section's two coordinates by
\(h\) changes its Frobenius coordinate to
\(U'=h^{-1}U\varphi_q(h)\).
On \(H_p\), the automorphism \(\varphi_q\)
is the identity, since its class-two action is
the pure cyclotomic action and \(q\equiv1\pmod p\).
The prescribed central image of \(U\) is
therefore also the image of \(U'\).
Equation (FB.6) in \(P\) now says
\[
 U'x^{qm}(U')^{-1}=x^{qm}.
 \tag{FB.9}
\]

For completeness, the needed centralizer
calculation is integral. Use the free
class-three Lie coordinates over \(\mathbb Z/M\),
with basis \(X_0,Y_0,C_0=[X_0,Y_0]\),
\(D_0=[C_0,X_0]\), \(E_0=[C_0,Y_0]\).
Write
\(\log U'=aX_0+bY_0+c_0C_0+d_0D_0+e_0E_0\).
The logarithm of the conjugate of \(x^{qm}\)
is
\[
 qmX_0-qmbC_0+
 qm\bigl(c_0+\tfrac12ab\bigr)D_0+
 \tfrac12qmb^2E_0.
 \tag{FB.10}
\]
Here \(v_p(qm)=k\). Equation (FB.9) first
forces \(b\equiv0\pmod p\), from the
\(C_0\) coordinate, and then
\(c_0\equiv0\pmod p\), from \(D_0\).
Its image in \(H_p\) consequently cannot
be a nonzero central element.
This contradicts the preserved Frobenius
image and proves the proposition.
\end{proof}

\begin{result}{C.2.16}{Proved}{A two-sided bound for the selected A.65 first failures}
\label{cor:ghost-selected-failure-bound}
For the selected pair of
Theorem~\ref{thm:ghost-unbounded-finite-failures},
let \(q\) be its distinguished first positive
support prime and \(m=E_0(q)\).
Its first empty unrestricted integral
nilpotent section fibre occurs at a class \(r\)
satisfying
\[
 R<r\le10\bigl(5^{\,v_5(m)+1}-1\bigr).
 \tag{FB.11}
\]
In the parameters of Theorem~\ref{thm:ghost-general-exponent-arithmetic},
put \(d=\max(R,3)-1\), \(n=J!\), and \(a=v_5(n)\).
The exact exponent is
\[
 v_5(m)=v_5(C)+da+v_5(d!).
 \tag{FB.12}
\]
\end{result}
\begin{proof}
Every integral class-two lift of the selected
pair has the nonzero central Frobenius image
at \(q\) proved in
Theorem~\ref{thm:ghost-unbounded-finite-failures}.
Its inertia abelian coordinates are \((m,0)\).
The actual projection from the full integral
class-\(B(5,m)\) arithmetic quotient to its
geometrically pro-five quotient, followed by
restriction to \(G_{\mathbb Q_q}\), would give
the section excluded by
Proposition~\ref{prop:ghost-finite-local-bound}.
Thus the former fibre is empty.
The same pair has a full integral class-\(R\)
section by Theorem~\ref{thm:ghost-fixed-depth-survival},
which proves (FB.11).
Finally the first-stage shell gives
\(v_5(q-1)=v_5(n)=a\).
The lifting-the-exponent identity gives
\(v_5(q^i-1)=a+v_5(i)\) for \(i\ge1\).
Substitution into
\(m=Cn\prod_{i=2}^d(q^i-1)\)
gives (FB.12).
\end{proof}
